\documentclass[11pt,letterpaper]{article}
\usepackage[margin=2.54cm]{geometry}
\usepackage{fontspec}
\setmainfont{Arial}
\usepackage[spanish,es-noshorthands]{babel}
\usepackage{xcolor}
\usepackage{booktabs}
\usepackage{tabularx}
\usepackage{colortbl}
\usepackage{graphicx}
\usepackage{hyperref}
\usepackage{microtype}
\usepackage{setspace}

\definecolor{unabblue}{HTML}{173864}
\definecolor{softgray}{HTML}{F9FAFC}
\hypersetup{colorlinks=true,linkcolor=unabblue,urlcolor=unabblue}

\onehalfspacing
\setlength{\parindent}{0pt}
\setlength{\parskip}{0.6em}

\begin{document}
\thispagestyle{empty}

\begin{center}
  \includegraphics[width=2.5cm,height=2.5cm,keepaspectratio]{/home/jaimeberrios/Proyectos/UNAB/Ciclo-6/sDSR43/Plantilla-Documentos-Investigacion/assets/img/logo-portada.png}\\[0.3cm]
  {\large\sffamily\bfseries\color{unabblue}UNIVERSIDAD DR. ANDRÉS BELLO}\\[0.1cm]
  {\normalsize\sffamily\bfseries\color{unabblue}FACULTAD DE TECNOLOGÍA E INNOVACIÓN}\\[0.1cm]
  {\small\sffamily\color{unabblue}DISEÑO DE SISTEMAS DE SEGURIDAD DE REDES DE DATOS}\\[0.4cm]
  {\Large\sffamily\bfseries\color{unabblue}BITÁCORA DE AVANCES Y MODIFICACIONES - FASE 2}\\[0.15cm]
  {\color{unabblue}\rule{0.6\textwidth}{1.2pt}}
\end{center}

\vspace{0.4cm}

\begin{table}[h!]
  \centering
  \small
  \begin{tabularx}{\textwidth}{@{}>{\bfseries\color{unabblue}}p{3.8cm} X@{}}
    \toprule
    Proyecto: & Diseño de un modelo de seguridad perimetral y segmentación lógica para el aseguramiento de la red de datos en la Sede de Enlaces con la Educación (Santa Rosa de Lima, 2026). \\
    \midrule
    Periodo / Fase: & Fase 2: Gestión de Riesgos y Políticas. \\
    \midrule
    Fecha de entrega: & 12 de octubre de 2026. \\
    \midrule
    Docente titular: & MduEV. Ing. Franklin Geovany Buruca Hernández. \\
    \midrule
    Equipo de trabajo: & Genesis Saravia, Jaime Berrios, Jefferson Bonilla, Rosmeri González, Florentín Treminio. \\
    \bottomrule
  \end{tabularx}
\end{table}

\vspace{0.3cm}

\section*{\color{unabblue}1. Propósito de la Bitácora}
El presente documento tiene como objetivo registrar de forma sistemática las correcciones, adaptaciones metodológicas y desarrollos técnicos ejecutados por el equipo investigador durante la Fase 2 del proyecto integrador de cátedra, atendiendo las directrices académicas y las observaciones emitidas por la cátedra en el primer entregable.

\section*{\color{unabblue}2. Registro de Actividades y Cambios Implementados}

\begin{table}[h!]
  \centering
  \small
  \renewcommand{\arraystretch}{1.3}
  \begin{tabularx}{\textwidth}{@{}>{\bfseries}p{3.2cm} X X@{}}
    \toprule
    \color{unabblue}\textbf{Aspecto / Módulo} & \color{unabblue}\textbf{Modificación o Avance Realizado} & \color{unabblue}\textbf{Impacto Metodológico} \\
    \midrule
    Objetivo Específico 4 & Sustitución del verbo «Implementar» por «Configurar controles de capa de aplicación... en un laboratorio simulado». & Aclara el alcance tecnológico sin comprometer la intervención física en las instalaciones de la sede. \\
    \midrule
    Carátula Institucional & Ajuste de la portada para presentar un único logotipo institucional UNAB centrado y debidamente ordenado. & Elimina duplicidad visual y estandariza la presentación según el manual de estilo institucional. \\
    \midrule
    Marco Metodológico (Sampieri) & Ratificación del diseño descriptivo y propositivo (no experimental y transversal) y justificación de no requerir hipótesis causal obligatoria. & Alineación estricta con las pautas de Hernández-Sampieri y la retroalimentación del docente. \\
    \midrule
    Análisis de Resultados (Fase 2) & Incorporación de la Matriz de Análisis y Evaluación de Riesgos formal categorizando 5 activos críticos y evaluando probabilidad e impacto. & Cumplimiento formal del entregable de la Fase 2 (Gestión de Riesgos) según la planificación académica. \\
    \midrule
    Instrumentos de Medición & Diseño y formalización del Cuestionario de Diagnóstico Técnico de Infraestructura y Seguridad de Red aplicado a informantes clave. & Respaldo metodológico documental para la recolección de datos y justificación técnica del diagnóstico. \\
    \midrule
    Anexos y Evidencias & Inclusión del registro fotográfico de las sesiones de trabajo del equipo (Anexos 1 y 2) y digitalización del cuestionario formal aplicado en campo a la informante clave (Anexo 3). & Constancia, trazabilidad y verificación fehaciente del trabajo colaborativo y de campo. \\
    \bottomrule
  \end{tabularx}
\end{table}

\vspace{0.3cm}

\section*{\color{unabblue}3. Estado Actual y Siguiente Fase}
Con la culminación de la Fase 2, el proyecto cuenta con un diagnóstico situacional consolidado, delimitación metodológica aprobada y una matriz de evaluación de riesgos formalizada. La siguiente etapa (Fase 3: Aseguramiento de Red y Transporte) comprenderá el diseño topológico en GNS3, la configuración de zonas lógicas (VLANs), direccionamiento IP y políticas de filtrado mediante cortafuegos GNU/Linux.

\end{document}
